% GENERATED FILE. Edit canonical lessons, then run npm run generate:books.
% canonical-source-hash: fnv1a64:e6569a2c4e58bf7b
% canonical-lessons: SA-C474-sira, SA-C474-snayuh, SA-C474-pesi, SA-C474-mamsam, SA-C474-spandate

\chapter{Veins and Muscles}
\label{ch:sa-veins-and-muscles}
\hlchaptermodality{474}

\begin{chapteropening}
\textbf{By the end of this chapter:} \emph{I can talk about blood, sinews and muscles.}

Complete the last lesson of chapter 474: I can talk about blood, sinews and muscles.
\end{chapteropening}

\section[\texorpdfstring{\sk{सिरा} (sirā)}{सिरा (sirā)}]{\sk{सिरा} (sirā) — a vein}

\label{lesson:SA-C474-sira}



\begin{quote}
Before the new one: say the Sanskrit for a fist, then the Sanskrit for bends a limb.
\end{quote}

\subsection*{You'll want to know: \sk{सिरा}}
\textbf{\sk{सिरा}} — \emph{sirā} — \textquotedblleft{}a vein\textquotedblright{}.

\begin{quote}
\sk{रक्तं} \sk{सिरासु} \sk{वहति।}
\end{quote}

\begin{quote}
\emph{raktaṃ sirāsu vahati.}
\end{quote}

\begin{quote}
Blood flows in the veins.
\end{quote}

A feminine noun.

\begin{cousinweb}
\textbf{\sk{सिरा}} (\emph{sirā}) comes from the root √sṛ, \textquotedblleft{}to flow\textquotedblright{}: MW's first sense is \textquotedblleft{}a stream\textquotedblright{}, then \textquotedblleft{}any tubular vessel of the body, a vein, an artery\textquotedblright{} (MW).
\end{cousinweb}

\begin{grammarlens}[title={What you've built}]
One more word for describing what you lived and felt.
\end{grammarlens}

\subsection*{Your turn}
\begin{itemize}
\raggedright
  \item \emph{Say it:} \emph{sirā}
  \item \emph{Say it:} \emph{sirā}, once more
  \item \emph{Say it:} \emph{raktaṃ sirāsu vahati.} — Blood flows in the veins.
  \item \emph{Recall:} say the Sanskrit for a fist, then the Sanskrit for bends a limb, then say \emph{sirā} again
\end{itemize}

\begin{tcolorbox}[breakable,colback=teal!4,colframe=teal!35!black,title={Before you move on}]
What does \sk{सिरा} mean? (\textquotedblleft{}A vein\textquotedblright{}.) Say it once more.
\end{tcolorbox}
\section[\texorpdfstring{\sk{स्नायुः} (snāyuḥ)}{स्नायुः (snāyuḥ)}]{\sk{स्नायुः} (snāyuḥ) — a sinew}

\label{lesson:SA-C474-snayuh}



\begin{quote}
Before the new one: say the Sanskrit for bends a limb, then the Sanskrit for a vein.
\end{quote}

\subsection*{You'll want to know: \sk{स्नायुः}}
\textbf{\sk{स्नायुः}} — \emph{snāyuḥ} — \textquotedblleft{}a sinew\textquotedblright{}.

\begin{quote}
\sk{स्नायुः} \sk{अस्थि} \sk{बध्नाति।}
\end{quote}

\begin{quote}
\emph{snāyuḥ asthi badhnāti.}
\end{quote}

\begin{quote}
Sinew binds the bone.
\end{quote}

A feminine noun.

\begin{cousinweb}
\textbf{\sk{स्नायुः}} (\emph{snāyuḥ}) MW compares, with a question mark, German \emph{Sehne} and English \emph{sinew}, and reports that some trace it to a root meaning \textquotedblleft{}to bind\textquotedblright{}; its origin is not settled (MW).
\end{cousinweb}

\begin{grammarlens}[title={What you've built}]
One more word for describing what you lived and felt.
\end{grammarlens}

\subsection*{Your turn}
\begin{itemize}
\raggedright
  \item \emph{Say it:} \emph{snāyuḥ}
  \item \emph{Say it:} \emph{snāyuḥ}, once more
  \item \emph{Say it:} \emph{snāyuḥ asthi badhnāti.} — Sinew binds the bone.
  \item \emph{Recall:} say the Sanskrit for bends a limb, then the Sanskrit for a vein, then say \emph{snāyuḥ} again
\end{itemize}

\begin{tcolorbox}[breakable,colback=teal!4,colframe=teal!35!black,title={Before you move on}]
What does \sk{स्नायुः} mean? (\textquotedblleft{}A sinew\textquotedblright{}.) Say it once more.
\end{tcolorbox}
\section[\texorpdfstring{\sk{पेशी} (peśī)}{पेशी (peśī)}]{\sk{पेशी} (peśī) — a muscle}

\label{lesson:SA-C474-pesi}



\begin{quote}
Before the new one: say the Sanskrit for a vein, then the Sanskrit for a sinew.
\end{quote}

\subsection*{You'll want to know: \sk{पेशी}}
\textbf{\sk{पेशी}} — \emph{peśī} — \textquotedblleft{}a muscle\textquotedblright{}.

\begin{quote}
\sk{व्यायामः} \sk{पेशीं} \sk{दृढां} \sk{करोति।}
\end{quote}

\begin{quote}
\emph{vyāyāmaḥ peśīṃ dṛḍhāṃ karoti.}
\end{quote}

\begin{quote}
Exercise makes the muscle firm.
\end{quote}

A feminine noun.

\begin{cousinweb}
\textbf{\sk{पेशी}} (\emph{peśī}) MW's first sense is \textquotedblleft{}a piece of flesh\textquotedblright{}, then \textquotedblleft{}a muscle\textquotedblright{}; MW gives no derivation, so its origin is uncertain (MW).
\end{cousinweb}

\begin{grammarlens}[title={What you've built}]
One more word for describing what you lived and felt.
\end{grammarlens}

\subsection*{Your turn}
\begin{itemize}
\raggedright
  \item \emph{Say it:} \emph{peśī}
  \item \emph{Say it:} \emph{peśī}, once more
  \item \emph{Say it:} \emph{vyāyāmaḥ peśīṃ dṛḍhāṃ karoti.} — Exercise makes the muscle firm.
  \item \emph{Recall:} say the Sanskrit for a vein, then the Sanskrit for a sinew, then say \emph{peśī} again
\end{itemize}

\begin{tcolorbox}[breakable,colback=teal!4,colframe=teal!35!black,title={Before you move on}]
What does \sk{पेशी} mean? (\textquotedblleft{}A muscle\textquotedblright{}.) Say it once more.
\end{tcolorbox}
\section[\texorpdfstring{\sk{मांसम्} (māṃsam)}{मांसम् (māṃsam)}]{\sk{मांसम्} (māṃsam) — flesh}

\label{lesson:SA-C474-mamsam}



\begin{quote}
Before the new one: say the Sanskrit for a sinew, then the Sanskrit for a muscle.
\end{quote}

\subsection*{You'll want to know: \sk{मांसम्}}
\textbf{\sk{मांसम्}} — \emph{māṃsam} — \textquotedblleft{}flesh\textquotedblright{}.

\begin{quote}
\sk{सिंहः} \sk{मांसं} \sk{खादति।}
\end{quote}

\begin{quote}
\emph{siṃhaḥ māṃsaṃ khādati.}
\end{quote}

\begin{quote}
A lion eats flesh.
\end{quote}

A neuter noun.

\begin{cousinweb}
\textbf{\sk{मांसम्}} (\emph{māṃsam}) MW gives \textquotedblleft{}flesh, meat\textquotedblright{}, from the Rigveda on; MW gives no derivation, so its origin is uncertain (MW).
\end{cousinweb}

\begin{grammarlens}[title={What you've built}]
One more word for describing what you lived and felt.
\end{grammarlens}

\subsection*{Your turn}
\begin{itemize}
\raggedright
  \item \emph{Say it:} \emph{māṃsam}
  \item \emph{Say it:} \emph{māṃsam}, once more
  \item \emph{Say it:} \emph{siṃhaḥ māṃsaṃ khādati.} — A lion eats flesh.
  \item \emph{Recall:} say the Sanskrit for a sinew, then the Sanskrit for a muscle, then say \emph{māṃsam} again
\end{itemize}

\begin{tcolorbox}[breakable,colback=teal!4,colframe=teal!35!black,title={Before you move on}]
What does \sk{मांसम्} mean? (\textquotedblleft{}Flesh\textquotedblright{}.) Say it once more.
\end{tcolorbox}
\section[\texorpdfstring{\sk{स्पन्दते} (spandate)}{स्पन्दते (spandate)}]{\sk{स्पन्दते} (spandate) — throbs}

\label{lesson:SA-C474-spandate}



\begin{quote}
Before the new one: say the Sanskrit for a muscle, then the Sanskrit for flesh.
\end{quote}

\subsection*{You'll want to know: \sk{स्पन्दते}}
\textbf{\sk{स्पन्दते}} — \emph{spandate} — \textquotedblleft{}throbs\textquotedblright{}.

\begin{quote}
\sk{भयेन} \sk{मम} \sk{हृदयं} \sk{स्पन्दते।}
\end{quote}

\begin{quote}
\emph{bhayena mama hṛdayaṃ spandate.}
\end{quote}

\begin{quote}
My heart throbs with fear.
\end{quote}

Past (it …): \sk{हृदयं} \sk{स्पन्दितवत्}, \emph{hṛdayaṃ spanditavat}.

\begin{cousinweb}
\textbf{\sk{स्पन्दते}} (\emph{spandate}) comes from the root √spand, \textquotedblleft{}to quiver, to throb\textquotedblright{} (MW).
\end{cousinweb}

\begin{grammarlens}[title={What you've built}]
One more word for describing what you lived and felt.
\end{grammarlens}

\subsection*{Your turn}
\begin{itemize}
\raggedright
  \item \emph{Say it:} \emph{spandate}
  \item \emph{Say it:} \emph{spandate}, once more
  \item \emph{Say it:} \emph{bhayena mama hṛdayaṃ spandate.} — My heart throbs with fear.
  \item \emph{Recall:} say the Sanskrit for a muscle, then the Sanskrit for flesh, then say \emph{spandate} again
\end{itemize}

\begin{tcolorbox}[breakable,colback=teal!4,colframe=teal!35!black,title={Before you move on}]
What does \sk{स्पन्दते} mean? (\textquotedblleft{}Throbs\textquotedblright{}.) Say it once more.
\end{tcolorbox}
